\documentclass[a4paper,11pt]{report}
\usepackage[utf8]{inputenc}
\usepackage{amsmath}
\usepackage{amsfonts,marvosym}
\usepackage{amssymb}
\usepackage{amsthm,pxfonts,textcomp}
\usepackage{mathrsfs} %permite o uso de letras trabalhadas
\usepackage[top=3cm,left=3cm,right=2cm,bottom=3cm]{geometry} %margens
\usepackage{caption}
\usepackage[brazil]{babel}
\usepackage{fontenc}
\usepackage{graphicx}

\usepackage{hyperref}
\begin{document}
 \thispagestyle{empty}
 \pagestyle{empty}%remove numeração de página
%Cabeçalho

\begin{center}
\begin{small}
\textbf
{Universidade Federal do Pará\\
Campus Universitário de Cametá\\
}\end{small}

\end{center}
%Caixa
\begin{flushleft}
\textbf{Professor (a): Elton Sarmanho}\\
\smallskip 
\textbf{Disciplina: Estrutura de dados I}\\
\begin{center}
 \textbf{Exercícios de ED I}
\end{center}
\end{flushleft}
\section{Exercícios de Análise de Complexidade}

\begin{enumerate}

\item Calcule a complexidade de tempo dos seguintes algoritmos e justifique sua resposta:

\begin{enumerate}
 \item[a)] Escreva um algoritmo para encontrar o maior elemento de um array de tamanho $n$. Qual é a sua complexidade?
 
 \item[b)] Escreva um algoritmo recursivo para encontrar o maior elemento de um array de tamanho $n$. Qual é a sua complexidade?
 
 \item[c)] Escreva um algoritmo para calcular $x^{y}$ (exponenciação). Qual é a sua complexidade em função de $y$?
 
 \item[d)] Escreva um algoritmo para verificar se uma palavra é um palíndromo. Qual é a sua complexidade?
\end{enumerate}

\end{enumerate}

\section{Exercícios de Recursividade}

\begin{enumerate}

\item Identifique o \textbf{caso base} e o \textbf{passo indutivo} nas seguintes definições recursivas:

\begin{enumerate}
 \item[a)] $n! = \begin{cases} 1, & \text{se } n = 0 \\ n \cdot (n-1)!, & \text{se } n > 0 \end{cases}$
 
 \item[b)] $f(n) = \begin{cases} 0, & \text{se } n = 0 \\ 1, & \text{se } n = 1 \\ f(n-1) + f(n-2), & \text{se } n > 1 \end{cases}$ (Fibonacci)
 
 \item[c)] Implemente uma função recursiva que encontre o \textbf{menor elemento} de um array de tamanho $n$. Defina claramente o caso base e o passo indutivo.
\end{enumerate}

\item Analise a seguinte implementação recursiva e determine sua complexidade de tempo:
\begin{enumerate}
 \item[a)] Uma função recursiva que calcula $n!$. Quantas chamadas recursivas são feitas para $n = 5$?
 
 \item[b)] Desenhe a \textbf{árvore de recurão} para o cálculo de $\text{Fibonacci}(4)$ e conte quantas vezes cada valor é calculado.
 \end{enumerate}

\item Implemente de forma recursiva:

\begin{enumerate}
 \item[a)] Um algoritmo que encontre todas as \textbf{substrings contíguas} de uma string de comprimento $n$. Quantas substrings existem?
 
 \item[b)] Uma função recursiva para gerar a sequência definida por:
  $$g_j = \begin{cases} j-1 & \text{se } 1 \leq j \leq k \\ g_{j-1} + g_{j-2} & \text{se } j > k \end{cases}$$
  onde $k$ é um parâmetro dado.
\end{enumerate}

\item Problemas com recursividade:

\begin{enumerate}
 \item[a)] O que é \textbf{Stack Overflow}? Em que circunstâncias uma função recursiva pode causar esse problema?
 
 \item[b)] Defina \textbf{tail recursion} (recursão de cauda) e explique como compiladores modernos a otimizam.
 
 \item[c)] Uma função recursiva sem caso base ou com um caso base incorreto pode causar que tipo de problema?
\end{enumerate}

\end{enumerate}

\section{Observações}


\end{document}